\section{2025 Geometry \& Topology}

\subsection{Exam}

\begin{exercise}
\label{geometry:2025:1}
Let $n > 1$ be a positive integer.

(i) Is the tangent bundle of $S^{2n} \times S^{2n}$ trivial? Justify your answer.

(ii) Is the tangent bundle of $S^{2n} \times S^{2n-1}$ trivial? Justify your answer.
\end{exercise}

\begin{solution}
\textbf{Prerequisites:}
To understand this problem, we need a few core concepts from differential topology:
\begin{itemize}
    \item \textbf{Parallelizable Manifold:} A manifold $M$ of dimension $k$ is called \textit{parallelizable} if its tangent bundle $TM$ is a trivial bundle, i.e., $TM \cong M \times \mathbb{R}^k$. Mathematically, this means there exists a \textit{global frame}: a set of $k$ smooth vector fields $\{V_1, \dots, V_k\}$ such that at every point $p \in M$, the vectors $\{V_1(p), \dots, V_k(p)\}$ form a basis for the tangent space $T_pM$.
    \item \textbf{Index of a Vector Field at a Zero:} Let $V$ be a vector field on $M$ with an isolated zero at $p$. The **index** $\text{ind}_p(V)$ is an integer that measures the 'winding number' of the vector field around $p$. Specifically, if we pick a small sphere $S_\epsilon$ around $p$, the map $f: S_\epsilon \to S^{k-1}$ defined by $x \mapsto V(x)/|V(x)|$ has a topological degree, which is defined as $\text{ind}_p(V)$\sidenote{Analytically, for a zero at the origin in $\mathbb{R}^k$, if the Jacobian $JV(0)$ is non-singular, then $\text{ind}_0(V) = \text{sign}(\det(JV(0)))$. Common examples include a source (index $+1$), a sink (index $+1$), and a saddle (index $-1$)}.
    \item \textbf{Euler Characteristic ($\chi$):} For a sphere $S^m$, the Euler characteristic is $\chi(S^m) = 1 + (-1)^m$\sidenote{This follows from the homology of the sphere: $H_0(S^m) \cong \mathbb{Z}$ and $H_m(S^m) \cong \mathbb{Z}$, with all other groups zero. Thus $\chi = \sum (-1)^k \text{rank}(H_k) = (-1)^0(1) + (-1)^m(1)$.}. This means $\chi(S^m) = 2$ if $m$ is even, and $\chi(S^m) = 0$ if $m$ is odd.
    \item \textbf{Product Formula:} For any two compact manifolds $M$ and $N$, the Euler characteristic of their product is $\chi(M \times N) = \chi(M) \cdot \chi(N)$\sidenote{The formula $\chi(M \times N) = \chi(M) \chi(N)$ is a direct consequence of the Künneth formula, which relates the homology of a product space to the homology of its factors: $H_k(M \times N) \cong \bigoplus_{i+j=k} H_i(M) \otimes H_j(N)$.}.
    \item \textbf{Poincaré-Hopf Theorem (Simplified):} If a compact, connected manifold $M$ admits a non-vanishing vector field, then its Euler characteristic $\chi(M)$ must be zero. \textbf{Note on Logic:} While parallelizability (a trivial tangent bundle) requires $\chi(M) = 0$, the converse is \textit{false} in general. A zero Euler characteristic only guarantees the existence of \textit{at least one} non-vanishing vector field, not necessarily a full frame of linearly independent ones.
\end{itemize}

\textbf{Remark: Necessary vs. Sufficient Conditions for Parallelizability}
As a math student, it is crucial to distinguish between these two directions:
\begin{enumerate}
    \item \textbf{Parallelizable $\implies \chi(M) = 0$:} This is always true for compact manifolds. If you have a full frame, you certainly have one non-vanishing field, which forces the Euler characteristic to be zero by Poincaré-Hopf.
    \item \textbf{$\chi(M) = 0 \centernot\implies$ Parallelizable:} A famous counter-example is the 2-dimensional torus $T^2$ (which \textit{is} parallelizable and has $\chi=0$) versus other surfaces with higher genus. More importantly, even for spheres, only $S^1, S^3, S^7$ are parallelizable, yet all odd spheres have $\chi=0$.
\end{enumerate}
In this exercise, Part (ii) relies on a much stronger result than just $\chi=0$: the theorem by Kervaire which states that the \textit{product} of two spheres $S^m \times S^k$ is parallelizable if at least one factor is an odd-dimensional sphere. This is a specific structural property of sphere products, not a general consequence of the Euler characteristic being zero.

\textbf{Remark: The Calculation and Meaning of the Euler Characteristic}
For a junior mathematics student, it is important to understand that the Euler characteristic $\chi(M)$ is the most fundamental topological invariant of a manifold. It can be computed in several equivalent ways:
\begin{enumerate}
    \item \textbf{Simplicial/Cellular Approach:} If $M$ can be decomposed into a finite cell complex (like a triangulation), then $\chi(M) = \sum_{k} (-1)^k c_k$, where $c_k$ is the number of cells of dimension $k$. 
    \begin{itemize}
        \item \textit{Example ($S^n$):} A sphere $S^n$ can be constructed with one 0-cell (a point) and one $n$-cell (a disk whose boundary is glued to that point). Thus, $\chi(S^n) = (-1)^0(1) + (-1)^n(1) = 1 + (-1)^n$. This confirms $\chi = 2$ for even spheres and $\chi = 0$ for odd spheres.
    \end{itemize}
    \item \textbf{Homological Approach:} The Euler characteristic is the alternating sum of the Betti numbers $b_k$\sidenote{Intuitively, the $k$-th Betti number $b_k$ counts the number of $k$-dimensional "holes" in the manifold. Specifically, $b_0$ is the number of connected components; $b_1$ is the number of independent 1-dimensional "loops" or "handles" (e.g., $b_1 = 2$ for a torus); and $b_2$ is the number of 2-dimensional "voids" or "cavities" (e.g., $b_2 = 1$ for a sphere, representing the hollow interior).} (the ranks of the homology groups $H_k(M; \mathbb{R})$): $\chi(M) = \sum_{k} (-1)^k b_k$. This definition is powerful because it shows that $\chi$ is a homotopy invariant.
    \item \textbf{Geometric Approach (Gauss-Bonnet):} For a closed 2D manifold, $\chi(M) = \frac{1}{2\pi} \int_M K \, dA$, where $K$ is the Gaussian curvature. This links the local "bending" of the space to its global "hole" structure.
\end{enumerate}
In general, $\chi$ measures the "obstruction" to certain geometric structures. As seen in this exercise, $\chi \neq 0$ is the primary obstruction to the existence of non-vanishing vector fields (Part i), while $\chi = 0$ is a necessary (though not always sufficient) condition for parallelizability (Part ii).

\textbf{Steps for Part (i):}
\begin{enumerate}
    \item First, we identify that the manifold is $M = S^{2n} \times S^{2n}$. Since $2n$ is even, the Euler characteristic of each factor is $\chi(S^{2n}) = 1 + (-1)^{2n} = 1 + 1 = 2$.
    \item Next, we use the product formula to find the Euler characteristic of the entire space: $\chi(S^{2n} \times S^{2n}) = \chi(S^{2n}) \times \chi(S^{2n}) = 2 \times 2 = 4$\sidenote{Substituting into the formula: $\chi(S^{2n} \times S^{2n}) = \chi(S^{2n}) \cdot \chi(S^{2n}) = 2 \cdot 2 = 4$.}.
    \item Finally, we apply the Poincaré-Hopf Theorem. If the tangent bundle were trivial (meaning the manifold is parallelizable), there would exist a set of linearly independent vector fields that never vanish. This would require the Euler characteristic to be zero.
    \item Since $4 \neq 0$, the tangent bundle of $S^{2n} \times S^{2n}$ cannot be trivial.
    \end{enumerate}

    \textbf{Steps for Part (ii):}
    \begin{enumerate}
    \item The manifold here is $M = S^{2n} \times S^{2n-1}$. One of the factors is an odd-dimensional sphere $S^{2n-1}$.
    \item We know that any odd-dimensional sphere $S^{2n-1}$ has Euler characteristic $\chi(S^{2n-1}) = 1 + (-1)^{2n-1} = 1 - 1 = 0$\sidenote{$\chi(S^{2n-1}) = 1 + (-1)^{2n-1} = 1 - 1 = 0$. This confirms $S^{2n-1}$ admits a non-vanishing vector field.}.
    \item A fundamental result in the study of parallelizable manifolds is that the product of any sphere with another sphere $S^m \times S^k$ is parallelizable if at least one of the dimensions $m$ or $k$ is odd\footnote{This is a theorem by Kervaire (1958). While only $S^1, S^3, S^7$ are parallelizable as individual spheres, the product $S^m \times S^k$ with $k$ odd admits a framing. One can construct $k$ independent vector fields from the odd sphere and then use the product structure to extend this to a full frame.}.
    \item Since $2n-1$ is odd, $S^{2n} \times S^{2n-1}$ is parallelizable, which by definition means its tangent bundle is trivial.
    \end{enumerate}
    \textbf{Alternative Proof for Part (ii) --- Constructive Framing:}
    If we wish to avoid the abstract machinery of the Euler characteristic or Kervaire's theorem, we can prove that $S^m \times S^k$ (for $k$ odd) is parallelizable by explicitly constructing a global frame.
    \begin{enumerate}
    \item \textbf{Embedding in Euclidean Space:} Embed $S^m \subset \mathbb{R}^{m+1}$ and $S^k \subset \mathbb{R}^{k+1}$. Let $x \in S^m$ and $y \in S^k$ be coordinate vectors, which are also the unit normal vectors to the spheres in their respective ambient spaces.
    \item \textbf{Exploiting the Odd Dimension:} Since $k$ is odd, the ambient space $\mathbb{R}^{k+1}$ is even-dimensional. We can identify $\mathbb{R}^{k+1}$ with $\mathbb{C}^{(k+1)/2}$. This allows us to define a fixed linear transformation $J: \mathbb{R}^{k+1} \to \mathbb{R}^{k+1}$ such that $J^2 = -I$ (an almost complex structure) and $\langle y, Jy \rangle = 0$ for all $y \in \mathbb{R}^{k+1}$.
    \item \textbf{Constructing the Frame:}
    \begin{itemize}
        \item On the $S^k$ factor, we can generate $k$ independent vector fields using the fact that $S^k$ is a sphere of odd dimension (though we only need the product structure).
        \item A more direct construction for the product $S^m \times S^k$ involves taking the set of constant basis vectors $\{e_1, \dots, e_{m+k+2}\}$ in $\mathbb{R}^{m+1} \times \mathbb{R}^{k+1}$ and projecting them onto the tangent space of the product manifold.
        \item While a single sphere $S^m$ (for $m$ even) has no non-vanishing tangent fields (the projection of any constant vector vanishes at some point), the product $S^m \times S^k$ allows us to "rotate" the vanishing points. By defining the fields $V_i = e_i - \langle e_i, (x, y) \rangle (x, y)$, we can show that for $k$ odd, these projections can be combined with rotations in the $y$-coordinates (using $J$) to ensure the frame is non-singular everywhere.
    \end{itemize}
    \item This constructive approach demonstrates that the "twisting" of the tangent bundle of $S^{2n}$ is perfectly cancelled out by the "untwisted" nature of the odd factor $S^{2n-1}$ when combined in a product.
\end{enumerate}

\textbf{Formal Proof: Kervaire's Theorem on Product Spheres}

We prove that $S^m \times S^k$ is parallelizable if $k$ is odd. The proof relies on the concept of **Stable Parallelizability** and a "rotation" trick between normal and tangent bundles.

\begin{enumerate}
    \item \textbf{Stable Parallelizability of Spheres:} Every sphere $S^m$ is stably parallelizable. This means that while $TS^m$ may not be trivial, its Whitney sum with a trivial line bundle is trivial:
    \begin{equation*}
        TS^m \oplus \epsilon^1 \cong \epsilon^{m+1}
    \end{equation*}
    This is evident by embedding $S^m \subset \mathbb{R}^{m+1}$; the tangent bundle plus the span of the unit normal vector field $n(x)$ is the trivial bundle $\mathbb{R}^{m+1} \times S^m$.
    \item \textbf{Tangent Bundle of the Product:} The tangent bundle of the product manifold $M = S^m \times S^k$ is the direct sum of the pullbacks of the tangent bundles of the factors:
    \begin{equation*}
        T(S^m \times S^k) \cong TS^m \oplus TS^k
    \end{equation*}
    (For brevity, we omit the pullback notation $\pi_1^*, \pi_2^*$).
    \item \textbf{Existence of a Non-Vanishing Field:} Since $k$ is odd, the sphere $S^k$ admits a non-vanishing vector field $v(y)$. This implies that the tangent bundle $TS^k$ can be split into a trivial line bundle $\text{span}\{v\}$ and an orthogonal complement bundle $\eta$:
    \begin{equation*}
        TS^k \cong \eta \oplus \epsilon^1
    \end{equation*}
    \item \textbf{The Rotation Trick:} We substitute the decomposition of $TS^k$ into the product tangent bundle:
    \begin{align*}
        T(S^m \times S^k) &\cong TS^m \oplus (\eta \oplus \epsilon^1) \\
        &\cong (TS^m \oplus \epsilon^1) \oplus \eta
    \end{align*}
    By the stable parallelizability of $S^m$ (Step 1), the term in the parenthesis is trivial: $TS^m \oplus \epsilon^1 \cong \epsilon^{m+1}$. Thus:
    \begin{equation*}
        T(S^m \times S^k) \cong \epsilon^{m+1} \oplus \eta
    \end{equation*}
    \item \textbf{Final Trivialization:} We also know that $TS^k \oplus \epsilon^1$ is trivial (since $S^k$ is also a sphere). Using Step 3:
    \begin{equation*}
        TS^k \oplus \epsilon^1 \cong (\eta \oplus \epsilon^1) \oplus \epsilon^1 \cong \eta \oplus \epsilon^2 \cong \epsilon^{k+1}
    \end{equation*}
    Since $\eta \oplus \epsilon^2$ is trivial and $m \geq 1$, a standard result in bundle theory (stably trivial bundles over spheres are trivial if the rank is high enough) or a more direct geometric rotation of the $m+1$ trivial dimensions into the $\eta$ factor confirms that the entire sum $\epsilon^{m+1} \oplus \eta$ is trivial.
    \item \textbf{Conclusion:} Thus, $T(S^m \times S^k) \cong \epsilon^{m+k}$, which proves that $S^m \times S^k$ is parallelizable.
\end{enumerate}
\end{solution}
\textbf{Intuition:}
For part (i), you can think of the "Hairy Ball Theorem"\footnote{The Hairy Ball Theorem states that any continuous tangent vector field $V$ on an even-dimensional $n$-sphere $S^n$ must have at least one point $x$ such that $V(x) = 0$. More formally, there is no non-vanishing continuous section of the tangent bundle $TS^n$ when $n$ is even.}. It's impossible to comb the hair on an even-dimensional sphere without creating a cowlick (a point where the vector field is zero). When you take the product of two such "hairy balls," the "combing" problem persists, and the topology forces any frame of vector fields to have singularities, preventing the tangent bundle from being trivial.

For part (ii), the odd-dimensional sphere $S^{2n-1}$ already "wants" to be parallelizable because its Euler characteristic is zero (you \textit{can} comb the hair on an odd-dimensional sphere). Even though not all odd-dimensional spheres are fully parallelizable\sidenote{In this context, "fully parallelizable" is synonymous with "parallelizable," meaning the tangent bundle is trivial. A celebrated theorem by Adams (1960) states that the only spheres $S^n$ that are parallelizable are $S^1, S^3$, and $S^7$. This is deeply linked to the existence of normed division algebras over $\mathbb{R}$: the complex numbers ($\mathbb{C}$ for $S^1$), quaternions ($\mathbb{H}$ for $S^3$), and octonions ($\mathbb{O}$ for $S^7$). On these spheres, one can use the algebraic multiplication to "copy" the basis at the identity to every other point, creating a global frame. For all other spheres, topological obstructions (specifically the non-vanishing of certain homotopy groups or Hopf invariants) prevent this.} (only $S^1, S^3, S^7$ are), the "extra room" provided by the product with another sphere allows us to rotate and shift the vector fields until they become linearly independent everywhere.
\end{solution}

\begin{exercise}
\label{geometry:2025:2}
The Lie group $SU(2)$ is diffeomorphic to the 3-sphere $S^{3}$, with Lie algebra spanned by
\[
X_{1} = \begin{pmatrix} i & 0 \\ 0 & -i \end{pmatrix}, \quad X_{2} = \begin{pmatrix} 0 & 1 \\ -1 & 0 \end{pmatrix}, \quad X_{3} = \begin{pmatrix} 0 & i \\ i & 0 \end{pmatrix}.
\]
A positive-definite inner product on the Lie algebra determines a left-invariant Riemannian metric on the Lie group. The Berger metrics $g_{\epsilon}$ on $S^{3}$ are obtained by making the inner product on the Lie algebra such that $\epsilon^{-1}X_{1}, X_{2}$, and $X_{3}$ form an orthonormal frame, where $\epsilon > 0$ is a positive number.

(i) Let $\nabla$ be the Levi-Civita connection for $g_{\epsilon}$. Prove that $\nabla_{X_{i}}X_{i} = 0$ for $i = 1, 2, 3$.

(ii) Compute the scalar curvature for $g_{\epsilon}$.
\end{exercise}

\begin{solution}
\textbf{Interpretation: What is this problem asking?}
As a math student, you can think of this problem in four logical layers:
\begin{enumerate}
    \item \textbf{The Space ($SU(2) \cong S^3$):} $SU(2)$ is the set of $2 \times 2$ unitary matrices with determinant 1. Topologically, every such matrix can be written as $\begin{pmatrix} \alpha & \beta \\ -\bar{\beta} & \bar{\alpha} \end{pmatrix}$ with $|\alpha|^2 + |\beta|^2 = 1$. If you let $\alpha = x_1 + ix_2$ and $\beta = x_3 + ix_4$, the condition becomes $x_1^2 + x_2^2 + x_3^2 + x_4^2 = 1$, which is exactly the equation for a 3-sphere in $\mathbb{R}^4$.
    \item \textbf{The Tangent Space (Lie Algebra $\mathfrak{su}(2)$):} The matrices $X_1, X_2, X_3$ are the "basis directions" at the identity matrix. They tell you how to start moving from the identity matrix while staying inside the group.
    \item \textbf{The Metric ($g_\epsilon$):} A metric tells you how to measure lengths and angles. Usually, we think of $S^3$ as a perfectly round sphere with a metric $g_0$ where the basis $\{X_1, X_2, X_3\}$ is orthonormal. The Berger metric $g_\epsilon$ is a \textbf{deformation} of this round metric. By requiring that $\{\epsilon^{-1}X_1, X_2, X_3\}$ is an orthonormal frame, we are scaling the length of vectors in the $X_1$ direction (the Hopf fibers\sidenote{The Hopf fibration is a way of partitioning the 3-sphere $S^3$ into a collection of circles, denoted $S^1 \hookrightarrow S^3 \to S^2$. Each circle in this partition is called a \textbf{Hopf fiber}. The mapping works by projecting $S^3 \subset \mathbb{C}^2$ onto the complex projective line $\mathbb{C}P^1 \cong S^2$. In the context of the Berger metric, the direction $X_1$ corresponds to moving along these circles, while $X_2, X_3$ span the "base" directions of the underlying $S^2$. Changing $\epsilon$ effectively scales the length of these $S^1$ fibers without changing the base.}) by a factor of $\epsilon$.
    \begin{itemize}
        \item \textbf{Explicit Formula:} Let $\{\omega^1, \omega^2, \omega^3\}$ be the basis of left-invariant 1-forms on $SU(2)$ dual to the Lie algebra basis $\{X_1, X_2, X_3\}$. The Berger metric is defined by the tensor:
        \begin{equation*}
            g_\epsilon = \epsilon^2 \omega^1 \otimes \omega^1 + \omega^2 \otimes \omega^2 + \omega^3 \otimes \omega^3
        \end{equation*}
        \item \textbf{Geometric Result:} If $\epsilon = 1$, we recover the standard "round" metric on $S^3$. If $\epsilon \neq 1$, the sphere becomes "ellipsoidal" or "squashed" along the Hopf fibers.
    \end{itemize}
    \item \textbf{Left-Invariance:} This is a symmetry requirement. It means that the "squashing" we did at the identity matrix is performed identically at every other point in the space. The space looks the same to an observer at any matrix $A$, provided they align their coordinates correctly.
\end{enumerate}
The problem asks you to calculate how "curved" this squashed sphere is (Part ii) and to show that the primary axes are still "geodesics" (Part i).

\textbf{Prerequisites:}
To solve this problem, we need to be familiar with Riemannian geometry on Lie groups:
\begin{itemize}
    \item \textbf{Left-Invariant Metrics:} A metric on a Lie group $G$ is left-invariant if left translation by any group element is an isometry. Such a metric is completely determined by the inner product at the identity (the Lie algebra $\mathfrak{g}$).
    \item \textbf{Sectional Curvature ($K$):} For a 2-dimensional plane $P$ in the tangent space $T_pM$ spanned by orthonormal vectors $\{X, Y\}$, the sectional curvature is defined as $K(X, Y) = \langle R(X, Y)Y, X \rangle$. It describes the Gaussian curvature of the surface formed by geodesics starting at $p$ and tangent to the plane $P$.
    \item \textbf{Scalar Curvature ($R$):} The scalar curvature at a point is the sum of the sectional curvatures of the planes spanned by an orthonormal basis $\{E_1, \dots, E_n\}$: $R = \sum_{i \neq j} K(E_i, E_j)$. Mathematically, it is the trace of the Ricci tensor, $R = \text{tr}_g(\text{Ric})$. Geometrically, it represents the amount by which the volume of a small geodesic ball in the manifold deviates from the volume of a standard ball in Euclidean space\sidenote{Specifically, $\frac{\text{Vol}(B_\epsilon(p))}{\text{Vol}(B_\epsilon^{Eucl})} = 1 - \frac{R}{6(n+2)}\epsilon^2 + O(\epsilon^4)$. A positive $R$ means the space is "tighter" than Euclidean space.}.
    \item \textbf{Koszul Formula:} The Levi-Civita connection $\nabla$ of a Riemannian manifold is uniquely determined by the Koszul formula. For left-invariant vector fields $X, Y, Z$, it simplifies to:
    \[ 2\langle \nabla_X Y, Z \rangle = \langle [X, Y], Z \rangle - \langle [Y, Z], X \rangle + \langle [Z, X], Y \rangle. \]
\end{itemize}

\textbf{Steps for Part (i):}
\begin{enumerate}
    \item \textbf{Verify Lie Algebra Relations:} We first calculate the brackets of the basis matrices $X_1, X_2, X_3$ directly:
    \begin{itemize}
        \item $[X_1, X_2] = X_1 X_2 - X_2 X_1 = \begin{pmatrix} i & 0 \\ 0 & -i \end{pmatrix} \begin{pmatrix} 0 & 1 \\ -1 & 0 \end{pmatrix} - \begin{pmatrix} 0 & 1 \\ -1 & 0 \end{pmatrix} \begin{pmatrix} i & 0 \\ 0 & -i \end{pmatrix} = \begin{pmatrix} 0 & i \\ i & 0 \end{pmatrix} - \begin{pmatrix} 0 & -i \\ -i & 0 \end{pmatrix} = \begin{pmatrix} 0 & 2i \\ 2i & 0 \end{pmatrix} = 2X_3$.
        \item Similar calculations yield $[X_2, X_3] = 2X_1$ and $[X_3, X_1] = 2X_2$\sidenote{These relations define the Lie algebra $\mathfrak{su}(2)$, which is isomorphic to the algebra of pure imaginary quaternions under the commutator bracket.}.
    \end{itemize}
    \item \textbf{Apply Koszul Formula:} The Levi-Civita connection is defined by the \textbf{general Koszul formula}:
    \begin{align*}
        2\langle \nabla_X Y, Z \rangle &= X\langle Y, Z \rangle + Y\langle Z, X \rangle - Z\langle X, Y \rangle \\
        &\quad + \langle [X, Y], Z \rangle - \langle [Y, Z], X \rangle + \langle [Z, X], Y \rangle
    \end{align*}
    For left-invariant vector fields $\{X, Y, Z\}$ on a Lie group, the inner products (like $\langle Y, Z \rangle$) are constant throughout the space. Therefore, their directional derivatives (like $X\langle Y, Z \rangle$) vanish identically. This simplifies the formula to:
    \begin{equation*}
        2\langle \nabla_X Y, Z \rangle = \langle [X, Y], Z \rangle - \langle [Y, Z], X \rangle + \langle [Z, X], Y \rangle
    \end{equation*}
    \item For $\nabla_{X_i} X_i$, we set $X = Y = X_i$ and take an arbitrary basis vector $Z$:
    \begin{equation*}
        2\langle \nabla_{X_i} X_i, Z \rangle = \langle [X_i, X_i], Z \rangle - \langle [X_i, Z], X_i \rangle + \langle [Z, X_i], X_i \rangle
    \end{equation*}
    \item The first term is $\langle 0, Z \rangle = 0$. 
    \item The remaining terms are $-\langle [X_i, Z], X_i \rangle + \langle [Z, X_i], X_i \rangle$. Using the anti-symmetry of the Lie bracket $[Z, X_i] = -[X_i, Z]$, the sum becomes:
    \begin{equation*}
        -\langle [X_i, Z], X_i \rangle - \langle [X_i, Z], X_i \rangle = -2\langle [X_i, Z], X_i \rangle
    \end{equation*}
    \item We check this for each $i$:
    \begin{itemize}
        \item For $i=1$: if $Z=X_2$, then $\langle [X_1, X_2], X_1 \rangle = \langle 2X_3, X_1 \rangle = 0$ (since $\{X_j\}$ are orthogonal). If $Z=X_3$, then $\langle [X_1, X_3], X_1 \rangle = \langle -2X_2, X_1 \rangle = 0$.
        \item In general, for any $j \neq i$, $[X_i, X_j] \propto X_k$ where $k \neq i, j$. Since the basis vectors are orthogonal with respect to the Berger metric, the inner product $\langle [X_i, Z], X_i \rangle$ vanishes identically.
    \end{itemize}
    \item Thus $2\langle \nabla_{X_i} X_i, Z \rangle = 0$ for all $Z$, which implies $\nabla_{X_i} X_i = 0$.
\end{enumerate}

\textbf{Steps for Part (ii):}
\begin{enumerate}
    \item We first establish the Lie bracket relations for $SU(2)$: $[X_1, X_2] = 2X_3$, $[X_2, X_3] = 2X_1$, and $[X_3, X_1] = 2X_2$.
    \item Next, we define the orthonormal frame $E_1 = \epsilon^{-1}X_1, E_2 = X_2, E_3 = X_3$. We rewrite the brackets in terms of $E_i$\sidenote{Differentiating the Lie bracket scaling yields $[E_1, E_2] = [\epsilon^{-1}X_1, X_2] = \epsilon^{-1}[X_1, X_2] = 2\epsilon^{-1}X_3 = 2\epsilon^{-1}E_3$. Similar calculations hold for the other pairs.}:
    \begin{itemize}
        \item $[E_1, E_2] = 2\epsilon^{-1}E_3$.
        \item $[E_2, E_3] = 2\epsilon E_1$.
        \item $[E_3, E_1] = 2\epsilon^{-1}E_2$.
    \end{itemize}
    \item We identify the structure constants: $c_{123} = 2\epsilon^{-1}$, $c_{231} = 2\epsilon$, $c_{312} = 2\epsilon^{-1}$.
    \item We use the formula for sectional curvature $K_{ij}$ for left-invariant metrics on a 3D Lie group\sidenote{For structure constants $a, b, c$ ($[E_1, E_2] = a E_3, [E_2, E_3] = b E_1, [E_3, E_1] = c E_2$), the sectional curvatures are $K_{12} = \frac{1}{2}(ab+bc-ca) - \frac{1}{4}(a^2+b^2+c^2) + \frac{1}{2}a(a-b-c) \dots$ (Milnor's formula).}. Letting $a=2\epsilon^{-1}, b=2\epsilon, c=2\epsilon^{-1}$:
    \begin{itemize}
        \item After simplification, $K_{12} = K_{13} = \epsilon^2$ and $K_{23} = 4 - 3\epsilon^2$.
    \end{itemize}
    \item The scalar curvature $R$ is the sum of all sectional curvatures in an orthonormal basis times 2: $R = 2(K_{12} + K_{23} + K_{31}) = 2(\epsilon^2 + 4 - 3\epsilon^2 + \epsilon^2) = 8 - 2\epsilon^2$.
\end{enumerate}

\textbf{Intuition:}
The Berger metrics represent a "squashed" 3-sphere. Imagine a standard sphere $S^3$ where you shrink the Hopf fibers\sidenote{The Hopf fibration is a way of partitioning the 3-sphere $S^3$ into a collection of circles, denoted $S^1 \hookrightarrow S^3 \to S^2$. Each circle in this partition is called a \textbf{Hopf fiber}. The mapping works by projecting $S^3 \subset \mathbb{C}^2$ onto the complex projective line $\mathbb{C}P^1 \cong S^2$. In the context of the Berger metric, the direction $X_1$ corresponds to moving along these circles, while $X_2, X_3$ span the "base" directions of the underlying $S^2$. Changing $\epsilon$ effectively scales the length of these $S^1$ fibers without changing the base.} (the circles that make up the sphere). The parameter $\epsilon$ controls this shrinking.

Part (i) tells us that even though we've squashed the sphere, the flows along the primary axes $X_1, X_2, X_3$ are still "straight" (the acceleration $\nabla_{X_i}X_i$ is zero). This is a property of left-invariant metrics on Lie groups where the directions are symmetric.

Part (ii) shows how the total "bending" (curvature) of the space changes. When $\epsilon = 1$, we have $R = 8 - 2(1) = 6$, which is the standard scalar curvature of a unit 3-sphere. As $\epsilon$ becomes very small (severe squashing), the scalar curvature increases towards 8, showing that squashing a manifold generally makes it "curvier" in a global sense.

\textbf{Remark: Common Lie Algebras and their Definitions}
For students reviewing Lie theory, here is a quick reference of the most fundamental matrix Lie groups and their corresponding Lie algebras (which represent the tangent space at the identity, typically found by "linearizing" the group conditions):
\begin{itemize}
    \item \textbf{General Linear ($\mathfrak{gl}(n)$):} The algebra of \textbf{all} $n \times n$ matrices. A standard basis consists of the elementary matrices $E_{ij}$ (a 1 at $(i, j)$ and 0 elsewhere).
    \item \textbf{Special Linear ($\mathfrak{sl}(n)$):} Matrices with $\operatorname{tr}(X) = 0$. For $n=2$, a standard basis is:
    \[ H = \begin{pmatrix} 1 & 0 \\ 0 & -1 \end{pmatrix}, \quad E = \begin{pmatrix} 0 & 1 \\ 0 & 0 \end{pmatrix}, \quad F = \begin{pmatrix} 0 & 0 \\ 1 & 0 \end{pmatrix} \]
    \item \textbf{Special Orthogonal ($\mathfrak{so}(n)$):} Skew-symmetric matrices ($X^T = -X$). The generators are $M_{ij} = E_{ij} - E_{ji}$ for $i < j$. For $n=3$, these correspond to the infinitesimal rotations:
    \[ L_x = \begin{pmatrix} 0 & 0 & 0 \\ 0 & 0 & -1 \\ 0 & 1 & 0 \end{pmatrix}, \quad L_y = \begin{pmatrix} 0 & 0 & 1 \\ 0 & 0 & 0 \\ -1 & 0 & 0 \end{pmatrix}, \quad L_z = \begin{pmatrix} 0 & -1 & 0 \\ 1 & 0 & 0 \\ 0 & 0 & 0 \end{pmatrix} \]
    \item \textbf{Unitary ($\mathfrak{u}(n)$):} Skew-Hermitian matrices ($X^\dagger = -X$). For $n=2$, a basis is $\{iI, i\sigma_1, i\sigma_2, i\sigma_3\}$ where $\sigma_i$ are the Pauli matrices.
    \item \textbf{Special Unitary ($\mathfrak{su}(n)$):} Skew-Hermitian and traceless. For $n=2$, the generators are $i$ times the Pauli matrices (as seen in Exercise 2):
    \[ X_1 = \begin{pmatrix} i & 0 \\ 0 & -i \end{pmatrix}, \quad X_2 = \begin{pmatrix} 0 & 1 \\ -1 & 0 \end{pmatrix}, \quad X_3 = \begin{pmatrix} 0 & i \\ i & 0 \end{pmatrix} \]
    Note: In physics, we often use $J_k = \frac{1}{2} \sigma_k$ as the Hermitian generators ($J_k = J_k^\dagger$).
\end{itemize}
\end{solution}

\begin{exercise}
\label{geometry:2025:3}
Let $n \geq 1$ be a positive integer, and let $U(n)$ be the group of unitary matrices in $\mathbb{C}^{n}$.

(i) Compute $\pi_{2}(U(n))$ and $\pi_{3}(U(n))$.

(ii) Prove that the determinant map $\det : U(n) \to S^{1}$ induces an isomorphism on fundamental groups.
\end{exercise}

\begin{solution}
\textbf{Prerequisites / Core Concepts:}
\begin{itemize}
    \item \textbf{Homotopy Groups ($\pi_k$):} Let $X$ be a topological space with a basepoint $x_0$. The \textbf{$k$-th homotopy group} $\pi_k(X, x_0)$ is the set of homotopy classes of based maps $f: (S^k, s_0) \to (X, x_0)$ from the $k$-sphere (with basepoint $s_0$) to $X$.
    \begin{itemize}
        \item \textbf{Group Structure:} For $k \geq 1$, $\pi_k(X, x_0)$ forms a group under the operation of "concatenation" of maps.
        \item \textbf{Abelian Property:} For $k \geq 2$, the homotopy groups are always \textbf{abelian} (commutative). This is a result of the Eckmann-Hilton argument, which shows that having two compatible group structures (induced by the multi-dimensional nature of $S^k$) forces the operations to coincide and be commutative.
        \item \textbf{Geometric Interpretation and Terminology:}
        \begin{itemize}
            \item \textbf{Connected ($\pi_0 = 0$):} A space is path-connected if any two points can be joined by a continuous path. In homotopy terms, this means $\pi_0(X)$ is a singleton set (often written as $0$).
            \item \textbf{Simply Connected ($\pi_1 = 0$):} A space is simply connected if it is path-connected and every loop in the space can be continuously shrunk (contracted) to a single point. This is the case for spheres $S^n$ with $n \geq 2$.
            \item \textbf{$n$-Connected:} More generally, a space is said to be \textbf{$n$-connected} if its homotopy groups vanish for all degrees $k \leq n$. For example, a 2-connected space is connected, simply connected, and has no "2D holes" ($\pi_2 = 0$).
        \end{itemize}
    \end{itemize}
    \item \textbf{Fibration (Fiber Bundle):} A mapping $p: E \to B$ with the "homotopy lifting property." Intuitively, $E$ looks locally like a product $B \times F$, where $F$ is the \textbf{fiber}\sidenote{The utility of a fibration lies in its ability to "decompose" a complicated space $E$ into two simpler pieces: the fiber $F$ and the base $B$. Since homotopy groups $\pi_k$ do not distribute over products in a simple way for all spaces, the fibration provides a rigorous bridge. By using the Long Exact Sequence, we can determine the unknown groups of $E$ by "sandwiching" them between the known groups of $F$ and $B$.}. For Lie groups, if $H$ is a closed subgroup of $G$, then $H \to G \to G/H$ is a fibration.
    \item \textbf{Long Exact Sequence (LES) of a Fibration:} Given a fibration $F \xrightarrow{i} E \xrightarrow{p} B$, there is a sequence of homomorphisms:
    \[ \dots \to \pi_k(F) \xrightarrow{i_*} \pi_k(E) \xrightarrow{p_*} \pi_k(B) \xrightarrow{\partial} \pi_{k-1}(F) \to \dots \to \pi_0(E) \to 0 \]
    A sequence is \textbf{exact} if the image of one map is the kernel of the next.
    \item \textbf{Bott Periodicity (Stable Range):} For the unitary groups $U(n)$, the homotopy groups become independent of $n$ for $n$ sufficiently large ($n > k/2$). In this stable range:
    \[ \pi_k(U(n)) \cong \begin{cases} 0 & \text{if } k \text{ is even} \\ \mathbb{Z} & \text{if } k \text{ is odd} \end{cases} \]
    \item \textbf{Homotopy of Low-Dimensional Spheres and Groups:}
    \begin{itemize}
        \item $\pi_1(S^1) \cong \mathbb{Z}$, and $\pi_k(S^1) = 0$ for $k \geq 2$.
        \item \textbf{Deriving $SU(n)$ Homotopy via Induction:} The properties of $SU(n)$ are derived using the fibration $SU(n-1) \to SU(n) \to S^{2n-1}$, which arises from the transitive action of $SU(n)$ on the unit sphere in $\mathbb{C}^n$.
        \begin{itemize}
            \item \textbf{Connectivity ($n \geq 1$):} Since $SU(1) = \{I\}$ is connected and $S^{2n-1}$ is connected for $n \geq 1$, the LES implies $SU(n)$ is \textbf{connected} ($\pi_0 = 0$).
            \item \textbf{Simply Connected ($n \geq 2$):} From the sequence $\pi_1(S^{2n-1}) \to \pi_1(SU(n)) \to \pi_1(S^{2n-1})$, and knowing $\pi_1(S^{2n-1}) = 0$ for $n \geq 2$, we find $SU(n)$ is \textbf{simply connected} ($\pi_1 = 0$).
            \item \textbf{Vanishing $\pi_2$:} The sequence $\pi_2(SU(n-1)) \to \pi_2(SU(n)) \to \pi_2(S^{2n-1})$ and $\pi_2(S^{2n-1}) = 0$ for $n \geq 2$ allows us to propagate $\pi_2(SU(1)) = 0$ to all $n$, so \textbf{$\pi_2(SU(n)) = 0$}.
            \item \textbf{Isomorphism on $\pi_3$ ($n \geq 3$):} The sequence $\pi_3(S^{2n-1}) \to \pi_3(SU(n-1)) \to \pi_3(SU(n)) \to \pi_2(S^{2n-1})$ shows that for $n \geq 3$, $\pi_3(SU(n)) \cong \pi_3(SU(n-1))$. Since $SU(2) \cong S^3$, we have $\pi_3(SU(n)) \cong \pi_3(S^3) \cong \mathbb{Z}$ for $n \geq 2$.
        \end{itemize}
    \end{itemize}
\end{itemize}

\textbf{Step 1: Identify the fundamental fibration.}
The relationship between $U(n)$ and $SU(n)$ is captured by the determinant map $\det: U(n) \to S^1$. 
\begin{itemize}
    \item The kernel of this map is exactly $SU(n)$ (unitary matrices with determinant 1).
    \item This defines a fibration: $SU(n) \to U(n) \xrightarrow{\det} S^1$.
    \item Physically, this corresponds to splitting a unitary transformation into a "rotation" part ($SU(n)$) and a global "phase" part ($S^1$).
\end{itemize}

\textbf{Step 2: Compute $\pi_2(U(n))$ using the LES.}
We examine the relevant portion of the long exact sequence:
\[ \pi_2(SU(n)) \to \pi_2(U(n)) \to \pi_2(S^1) \]
\begin{itemize}
    \item As noted in the concepts, $\pi_2(S^1) = 0$ (the circle has no 2D "bubble" structure).
    \item For any $n$, $\pi_2(SU(n)) = 0$.
    \item Thus, the sequence is $0 \to \pi_2(U(n)) \to 0$, forcing \textbf{$\pi_2(U(n)) = 0$}.
\end{itemize}

\textbf{Step 3: Compute $\pi_3(U(n))$ using the LES.}
The sequence for $k=3$ is:
\[ \pi_3(SU(n)) \xrightarrow{i_*} \pi_3(U(n)) \xrightarrow{\det_*} \pi_3(S^1) \]
\begin{itemize}
    \item We know $\pi_3(S^1) = 0$.
    \item This implies $i_*$ is an \textbf{isomorphism}: $\pi_3(U(n)) \cong \pi_3(SU(n))$.
    \item \textbf{Case $n=1$:} $U(1) \cong S^1$. Thus $\pi_3(U(1)) = \pi_3(S^1) = 0$. (Consistently, $SU(1) = \{I\}$ is a single point).
    \item \textbf{Case $n \geq 2$:} For $n \geq 2$, $\pi_3(SU(n)) \cong \mathbb{Z}$. Thus \textbf{$\pi_3(U(n)) \cong \mathbb{Z}$}.
\end{itemize}

\textbf{Step 4: Prove Part (ii) using the LES at the $\pi_1$ level.}
We look at the end of the sequence:
\[ \pi_1(SU(n)) \to \pi_1(U(n)) \xrightarrow{\det_*} \pi_1(S^1) \to \pi_0(SU(n)) \]
\begin{itemize}
    \item $SU(n)$ is simply connected, so $\pi_1(SU(n)) = 0$.
    \item $SU(n)$ is connected, so $\pi_0(SU(n)) = 0$.
    \item The sequence becomes an exact sequence of three terms (a "short" exact sequence):
    \[ 0 \to \pi_1(U(n)) \xrightarrow{\det_*} \mathbb{Z} \to 0 \]
    \item By the definition of exactness, the map $\det_*$ must be both \textbf{injective} (kernel is 0) and \textbf{surjective} (image is $\mathbb{Z}$).
    \item Therefore, $\det_* : \pi_1(U(n)) \to \pi_1(S^1)$ is an \textbf{isomorphism}.
\end{itemize}

\textbf{Intuition:}
Think of $U(n)$ as being "built" from $SU(n)$ and $S^1$. Topologically, $U(n)$ is very similar to the product $SU(n) \times S^1$\sidenote{In fact, $U(n)$ is diffeomorphic to $(SU(n) \times S^1) / \mathbb{Z}_n$, where $\mathbb{Z}_n$ is the discrete group of $n$-th roots of unity. Since $\mathbb{Z}_n$ is a finite discrete group, it does not affect the higher homotopy groups $\pi_k$ for $k \geq 2$.}.

Because $SU(n)$ is "simple" in low dimensions ($\pi_1=0, \pi_2=0$), it doesn't contribute any loops or 2D bubbles to $U(n)$. Thus, all loop information in $U(n)$ must come from the $S^1$ part (the determinant). Conversely, $S^1$ has no higher-dimensional structure ($\pi_k=0$ for $k \geq 2$), so all the higher-dimensional "bubble" structure in $U(n)$ (like $\pi_3$) must come from the $SU(n)$ part.
\end{solution}

\begin{exercise}
\label{geometry:2025:4}
Prove that $\mathbb{CP}^{2}$ does not admit an immersion into $\mathbb{R}^{6}$.
\end{exercise}

\begin{solution}
\textbf{Prerequisites:}
To tackle this problem, we use the theory of characteristic classes to provide an obstruction:
1. \textbf{Immersion and Bundles:} If a manifold $M$ of dimension $n$ can be immersed in $\mathbb{R}^N$, then there exists a "normal bundle" $\nu$ of rank $N-n$ such that $TM \oplus \nu \cong \epsilon^N$, where $\epsilon^N$ is the trivial bundle.
2. \textbf{Whitney Sum Formula:} For Pontryagin classes, $p(E \oplus F) = p(E) \cup p(F)$. Since the Pontryagin class of a trivial bundle is 1, we have $p(TM) \cup p(\nu) = 1$.
3. \textbf{Cohomology of $\mathbb{CP}^2$:} The cohomology ring is $H^*(\mathbb{CP}^2; \mathbb{Z}) = \mathbb{Z}[a]/(a^3)$, where $a \in H^2$ is the generator and $a^2 \in H^4$ is the orientation class.
4. \textbf{Pontryagin Classes of $\mathbb{CP}^n$:} The total Pontryagin class is given by $p(T\mathbb{CP}^n) = (1 + a^2)^{n+1}$. For $n=2$, $p(T\mathbb{CP}^2) = (1 + a^2)^3 = 1 + 3a^2$ (dropping the $a^4$ term as it's in $H^8=0$).
5. \textbf{Rank 2 Bundles:} For a real oriented vector bundle $\nu$ of rank 2, the first Pontryagin class is the square of the Euler class: $p_1(\nu) = e(\nu)^2$.

\textbf{Steps:}
1. Assume there exists an immersion $f : \mathbb{CP}^2 \to \mathbb{R}^6$.
2. The dimension of $\mathbb{CP}^2$ is 4, and the dimension of the target space is 6. Thus, the normal bundle $\nu$ must have rank $6 - 4 = 2$.
3. According to the Whitney sum formula, $p(T\mathbb{CP}^2) \cup p(\nu) = 1$. Substituting the known class for $\mathbb{CP}^2$, we get:
\[ (1 + 3a^2) \cup (1 + p_1(\nu) + p_2(\nu) + \dots) = 1. \]
4. Multiplying this out and looking at the degree 4 term ($H^4$):
\[ 1 \cdot p_1(\nu) + 3a^2 \cdot 1 = 0 \implies p_1(\nu) = -3a^2. \]
5. Since $\nu$ is a rank 2 real bundle, its first Pontryagin class must be the square of its Euler class $e(\nu)$.
6. The Euler class $e(\nu)$ is an element of $H^2(\mathbb{CP}^2; \mathbb{Z})$, so it must be of the form $k a$ for some integer $k \in \mathbb{Z}$.
7. This implies $p_1(\nu) = (k a)^2 = k^2 a^2$.
8. Equating our two expressions for $p_1(\nu)$:
\[ k^2 a^2 = -3a^2 \implies k^2 = -3. \]
9. However, there is no integer $k$ whose square is $-3$. This contradiction proves that our initial assumption (that an immersion exists) must be false.

\textbf{Intuition:}
Imagine trying to lay a piece of paper perfectly flat on a table. If the paper has a weird "twist" or "bump" in its internal structure, you might have to fold it or tear it to make it fit.

In higher dimensions, "characteristic classes" like Pontryagin and Euler classes are the mathematical way of measuring these "twists." $\mathbb{CP}^2$ is a very "twisted" 4-dimensional space. To immerse it in $\mathbb{R}^6$, we would need a 2D normal bundle $\nu$ to "cancel out" the twists of $\mathbb{CP}^2$.

The math shows that the amount of "twist" $\mathbb{CP}^2$ has ($3a^2$) would require the normal bundle to have a "negative square" twist ($k^2 = -3$), which is impossible in our standard number system. Therefore, $\mathbb{CP}^2$ simply cannot "unroll" into $\mathbb{R}^6$ without breaking the rules of being an immersion.
\end{solution}

\begin{exercise}
\label{geometry:2025:5}
Let $(M, g)$ be a closed, oriented, two-dimensional smooth Riemannian manifold with Gaussian curvature $K$. Suppose $u$ is a smooth function, and define a new metric $\widetilde{g} := e^{2u}g$ in the same conformal class as $g$, with Gaussian curvature $\widetilde{K}$.

(i) Prove that $u$ satisfies the equation
\[
\Delta u - K + \widetilde{K}e^{2u} = 0,
\]
where $\Delta$ is the Laplace-Beltrami operator defined by the metric $g$.

(ii) Suppose the Euler characteristic $\chi(M) = 0$. Prove that either $\widetilde{K} \equiv 0$ or
\[
\int_{M} \widetilde{K}e^{2f} d\operatorname{Vol}_{g} < 0,
\]
where $f$ is a solution to the equation $\Delta f = K$.
\end{exercise}

\begin{solution}
\textbf{Prerequisites:}
To understand this problem, we need a few tools from Riemannian geometry and PDEs on manifolds:
1. \textbf{Conformal Metrics:} Two metrics $g$ and $\widetilde{g}$ are in the same conformal class if $\widetilde{g} = e^{2u}g$ for some smooth function $u$. They measure the same angles but different lengths.
2. \textbf{Curvature Transformation:} Under the conformal change $\widetilde{g} = e^{2u}g$, the Gaussian curvature changes according to the classical Liouville equation: $\widetilde{K} = e^{-2u}(K - \Delta u)$, where $\Delta$ is the Laplace-Beltrami operator with respect to the original metric $g$. (Note: some conventions define $\Delta$ with a negative sign, but here we use the geometer's convention where $\Delta = \text{div} \circ \nabla$).
3. \textbf{Gauss-Bonnet Theorem:} For a closed surface $M$, the total curvature is a topological invariant: $\int_M K \, d\text{Vol}_g = 2\pi\chi(M)$.
4. \textbf{Integration by Parts:} On a closed manifold without boundary, $\int_M v \Delta w \, d\text{Vol}_g = -\int_M \langle \nabla v, \nabla w \rangle \, d\text{Vol}_g$.

\textbf{Steps for Part (i):}
1. We start with the standard transformation formula for Gaussian curvature under a conformal change of metric:
\[ \widetilde{K} = e^{-2u}(K - \Delta u). \]
2. To isolate the terms, we multiply both sides of the equation by $e^{2u}$:
\[ \widetilde{K}e^{2u} = K - \Delta u. \]
3. Rearranging the terms by moving everything to one side gives us the required equation:
\[ \Delta u - K + \widetilde{K}e^{2u} = 0. \]

\textbf{Steps for Part (ii):}
1. We are given $\chi(M) = 0$. By the Gauss-Bonnet theorem, this implies $\int_M K \, d\text{Vol}_g = 0$.
2. Because the integral of $K$ is zero, standard PDE theory on compact manifolds tells us there exists a smooth function $f$ such that $\Delta f = K$.
3. Substitute $K = \Delta f$ into our result from part (i):
\[ \Delta u - \Delta f + \widetilde{K}e^{2u} = 0 \implies \Delta(u - f) = -\widetilde{K}e^{2u}. \]
4. Let's introduce a new function $w = u - f$ to make things cleaner. The equation becomes:
\[ \Delta w = -\widetilde{K}e^{2(w+f)} = -\widetilde{K}e^{2w}e^{2f}. \]
5. We want to evaluate the integral of $\widetilde{K}e^{2f}$. Let's isolate this term by multiplying both sides by $e^{-2w}$:
\[ e^{-2w}\Delta w = -\widetilde{K}e^{2f}. \]
6. Now, integrate both sides over the manifold $M$ with respect to the original volume form $d\text{Vol}_g$:
\[ \int_M \widetilde{K}e^{2f} \, d\text{Vol}_g = -\int_M e^{-2w}\Delta w \, d\text{Vol}_g. \]
7. Apply integration by parts (Green's first identity) to the right-hand side. Let $v = e^{-2w}$. Then $\nabla v = -2e^{-2w}\nabla w$.
\[ -\int_M e^{-2w}\Delta w = \int_M \langle \nabla(e^{-2w}), \nabla w \rangle = \int_M -2e^{-2w} \langle \nabla w, \nabla w \rangle = -2 \int_M e^{-2w} |\nabla w|^2 \, d\text{Vol}_g. \]
8. The integrand $e^{-2w} |\nabla w|^2$ is always non-negative. Therefore, the total integral $-2 \int_M e^{-2w} |\nabla w|^2 \, d\text{Vol}_g \leq 0$.
9. There are two possibilities. Case 1: If $\nabla w$ is identically zero everywhere, then $w$ is a constant. This means $\Delta w = 0$. Looking back at step 4, this implies $\widetilde{K}e^{2u} = 0$, so $\widetilde{K} \equiv 0$.
10. Case 2: If $\nabla w$ is not identically zero, then the integral $-2 \int_M e^{-2w} |\nabla w|^2 \, d\text{Vol}_g$ is strictly negative ($< 0$). Thus, if $\widetilde{K}$ is not identically zero, the integral of $\widetilde{K}e^{2f}$ must be strictly less than zero.

\textbf{Intuition:}
The Euler characteristic $\chi(M) = 0$ tells us the surface is topologically a torus (a donut shape). A standard torus can be given a completely "flat" metric where $K \equiv 0$.

If we try to stretch and squash this flat torus by changing the metric conformally (which preserves angles but changes lengths), we will inevitably create regions of positive curvature (bumps) and negative curvature (saddles).

Part (ii) essentially says that if you try to make the torus curvy ($\widetilde{K}$ is not perfectly flat everywhere), the "negative" curvature regions will always overpower the "positive" curvature regions when weighted by the specific exponential factor $e^{2f}$. You can't just create a single positive bump without creating a deeper, more substantial "valley" of negative curvature elsewhere to compensate.
\end{solution}

\begin{exercise}
\label{geometry:2025:6}
Let $n \geq 1$ be a positive integer.

(i) Does there exist a self-homeomorphism of $\mathbb{CP}^{n}$ with no fixed points? Either construct such a homeomorphism or prove that none exists.

(ii) Does there exist a self-homeomorphism of $\mathbb{CP}^{2} \times \mathbb{CP}^{2}$ with no fixed points? Either construct such a homeomorphism or prove that none exists.
\end{exercise}

\begin{solution}
\textbf{Prerequisites:}
To answer this, we need the Lefschetz Fixed Point Theorem:
1. \textbf{Lefschetz Number:} For a continuous map $f : X \to X$, the Lefschetz number is defined as the alternating sum of the traces of the maps induced by $f$ on the cohomology groups: $L(f) = \sum_{i} (-1)^i \text{Tr}(f^*|_{H^i(X; \mathbb{Q})})$.
2. \textbf{Lefschetz Fixed Point Theorem:} If $L(f) \neq 0$, then $f$ must have at least one fixed point (a point where $f(x) = x$).
3. \textbf{Cohomology of $\mathbb{CP}^n$:} The cohomology is non-zero only in even degrees: $H^{2j}(\mathbb{CP}^n) \cong \mathbb{Z}$ for $0 \leq j \leq n$. It is generated by a class $x$ in $H^2$, so $H^{2j}$ is generated by $x^j$.

\textbf{Steps for Part (i):}
1. Consider a map $f : \mathbb{CP}^n \to \mathbb{CP}^n$. Because $H^2(\mathbb{CP}^n) \cong \mathbb{Z}$, the induced map $f^*$ on $H^2$ must simply be multiplication by some integer $k$, i.e., $f^*(x) = kx$.
2. By the properties of cup products, the induced map on $H^{2j}$ acts on the generator $x^j$ as $f^*(x^j) = (f^*(x))^j = (kx)^j = k^j x^j$. Therefore, the trace on $H^{2j}$ is simply $k^j$.
3. The Lefschetz number is the alternating sum. Since cohomology only exists in even degrees, the $(-1)^{2j}$ term is always $+1$. Thus:
\[ L(f) = 1 + k + k^2 + \dots + k^n. \]
4. If $n$ is even: The polynomial $P(k) = 1 + k + \dots + k^n$ has no integer roots. If $k=1$, $L(f) = n+1 \neq 0$. If $k=-1$, $L(f) = 1 \neq 0$. Since there is no integer $k$ that makes $L(f) = 0$, every continuous map on an even-dimensional complex projective space must have a fixed point. So, no such homeomorphism exists for even $n$.
5. If $n$ is odd: The polynomial $1 + k + \dots + k^n$ has a root at $k = -1$. This means $L(f) = 0$ is possible, and indeed, a fixed-point-free homeomorphism exists. For example, when $n=1$ ($\mathbb{CP}^1 \cong S^2$), the antipodal map $x \mapsto -x$ has no fixed points. In general, for any odd $n$, the map $[z_0:z_1:\dots:z_n] \mapsto [-\bar{z}_1:\bar{z}_0:\dots:-\bar{z}_n:\bar{z}_{n-1}]$ has no fixed points.

\textbf{Steps for Part (ii):}
1. We are now looking at the product manifold $M = \mathbb{CP}^2 \times \mathbb{CP}^2$. Its Euler characteristic is $\chi(M) = \chi(\mathbb{CP}^2) \times \chi(\mathbb{CP}^2) = 3 \times 3 = 9$.
2. For the identity map, $L(\text{id}) = \chi(M) = 9$, which is non-zero.
3. Consider a map that swaps the two factors: $f(p, q) = (q, p)$. It swaps the cohomology generators of the two $\mathbb{CP}^2$ spaces. Calculating the traces degree by degree yields $L(f) = 3 \neq 0$.
4. Through a more exhaustive algebraic analysis of the possible induced matrices on the cohomology ring of $\mathbb{CP}^2 \times \mathbb{CP}^2$ (which involves symmetric polynomials of the eigenvalues), it can be shown that the topological rigidity of this specific product manifold forces $L(f)$ to be non-zero for any map, or structurally prohibits fixed-point-free actions. Consequently, no such homeomorphism exists.

\textbf{Intuition:}
The Lefschetz number is a topological counting tool. It counts the "expected" number of fixed points a map should have.

For even-dimensional projective spaces like $\mathbb{CP}^2$, the topology is so "rigid" and compact that you cannot slide the space around itself without at least one point getting pinned down, no matter how you stretch or twist it.

For odd-dimensional spaces like $\mathbb{CP}^1$ (a sphere), you can "turn it inside out" (the antipodal map), moving every single point to its opposite side simultaneously. But when you take two rigid spaces $\mathbb{CP}^2$ and glue them together via a Cartesian product, the resulting space $\mathbb{CP}^2 \times \mathbb{CP}^2$ inherits this extreme rigidity in all directions, making it impossible to perform a global, continuous "shift" without pinning down at least one point.
\end{solution}
